\documentclass{article}
\usepackage[margin=1in, a4paper]{geometry}
\usepackage{amsmath, amssymb, amsfonts}
\usepackage{graphicx}
\usepackage{xcolor}
\usepackage{listings}
\usepackage{booktabs}
\usepackage[utf8]{inputenc}
\usepackage[T1]{fontenc}
\usepackage{hyperref}
\hypersetup{colorlinks=true, urlcolor=blue, linkcolor=blue}
\usepackage{enumitem}
\usepackage{float}

\definecolor{codegreen}{rgb}{0,0.6,0}
\definecolor{codegray}{rgb}{0.5,0.5,0.5}
\definecolor{codepurple}{rgb}{0.58,0,0.82}
\definecolor{backcolour}{rgb}{0.99,0.99,0.99}
% Professional color palette for academic publishing
\definecolor{myblue}{cmyk}{0.8,0.4,0,0.1}
\definecolor{mygreen}{cmyk}{0.7,0,0.5,0.2}
\definecolor{myorange}{cmyk}{0,0.5,0.8,0.1}
\definecolor{mygray}{cmyk}{0,0,0,0.4}
\definecolor{arrowcolor}{cmyk}{0.9,0.5,0,0.3}
\definecolor{nodecolor}{cmyk}{0.6,0.2,0,0.05}
\definecolor{framecolor}{cmyk}{0.4,0.1,0,0.1}

\lstdefinestyle{mystyle}{
    backgroundcolor=\color{backcolour},   
    commentstyle=\color{codegreen},
    keywordstyle=\color{magenta},
    numberstyle=\tiny\color{codegray},
    stringstyle=\color{codepurple},
    basicstyle=\ttfamily\footnotesize,
    breakatwhitespace=false,         
    breaklines=true,                 
    captionpos=b,                    
    keepspaces=true,                 
    numbers=left,                    
    numbersep=5pt,                  
    showspaces=false,                
    showstringspaces=false,
    showtabs=false,                  
    tabsize=2
}
\lstset{style=mystyle}

\title{The Classic Economic Order Quantity (EOQ) Problem}
\author{The Logistics \& Supply Chain Problem-Solving Playbook}
\date{}

\begin{document}

\maketitle

\section{The Classic Economic Order Quantity (EOQ) Problem}

The Economic Order Quantity (EOQ) problem represents one of the most fundamental challenges in inventory management, serving as the cornerstone of quantitative supply chain optimization. This classic problem seeks to determine the optimal order quantity that minimizes the total cost of inventory management by balancing two competing cost components: ordering costs and holding costs. The EOQ model has been the foundation of inventory theory since its development in 1913 by Ford W. Harris, and continues to be relevant in modern supply chain management across industries ranging from retail to manufacturing.

\subsection{The Scenario: Balancing the Inventory Cost Equation}

Consider MegaMart, a large retail chain that sells premium coffee beans to consumers nationwide. The company faces a critical inventory management challenge for one of its bestselling products: Colombian Supremo coffee beans. MegaMart purchases these beans from a single supplier in Colombia and must decide how much to order each time to minimize total inventory costs.

The company's supply chain team has gathered the following operational data: MegaMart sells approximately 12,000 bags of Colombian Supremo coffee per year, with demand remaining relatively stable throughout the year. Each time the company places an order with the Colombian supplier, it incurs fixed costs of \$150, which includes administrative processing, international shipping coordination, customs handling, and quality inspection procedures. Additionally, MegaMart incurs holding costs of \$8 per bag per year, encompassing warehouse storage fees, insurance, security, spoilage risk, and the opportunity cost of capital tied up in inventory.

The fundamental question facing MegaMart's supply chain team is: What is the optimal order quantity that minimizes the total annual cost of managing this inventory? This scenario exemplifies the classic trade-off in inventory management where ordering too frequently increases ordering costs, while ordering too infrequently increases holding costs due to excess inventory storage.

\subsection{Tier 1: The Pen \& Paper Method (Mathematical Formulation)}

The Economic Order Quantity problem can be elegantly formulated as a continuous optimization model that seeks to minimize the total annual cost function. This mathematical approach provides the theoretical foundation for understanding the optimal balance between ordering and holding costs.

\textbf{Sets and Indices:}
\begin{align}
T &= \text{Planning horizon (typically one year)}
\end{align}

\textbf{Parameters:}
\begin{align}
D &= \text{Annual demand rate (units per year)} \\
S &= \text{Ordering cost per order (\$ per order)} \\
H &= \text{Holding cost per unit per year (\$/unit/year)} \\
c &= \text{Unit cost (\$/unit)}
\end{align}

\textbf{Decision Variables:}
\begin{align}
Q &= \text{Order quantity (units per order)} \\
n &= \text{Number of orders per year} = \frac{D}{Q}
\end{align}

\textbf{Objective Function:}
The total annual cost function consists of three components: ordering costs, holding costs, and purchasing costs.

\begin{align}
\text{Minimize: } TC(Q) &= \text{Ordering Cost} + \text{Holding Cost} + \text{Purchasing Cost}
\end{align}

\begin{align}
TC(Q) &= \frac{D}{Q} \cdot S + \frac{Q}{2} \cdot H + D \cdot c
\end{align}

Where:
- \(\frac{D}{Q} \cdot S\) represents the annual ordering cost
- \(\frac{Q}{2} \cdot H\) represents the annual holding cost (assuming constant demand and linear inventory depletion)
- \(D \cdot c\) represents the annual purchasing cost (constant for given demand)

Since the purchasing cost \(D \cdot c\) is independent of the order quantity \(Q\), the optimization problem reduces to:

\begin{align}
\text{Minimize: } TC'(Q) &= \frac{D \cdot S}{Q} + \frac{Q \cdot H}{2}
\end{align}

\textbf{Constraints:}
\begin{align}
Q &\geq 0 \quad \text{(Non-negativity constraint)}
\end{align}

\textbf{Optimality Conditions:}
To find the optimal order quantity, we take the first derivative of the total cost function with respect to \(Q\) and set it equal to zero:

\begin{align}
\frac{dTC'(Q)}{dQ} &= -\frac{D \cdot S}{Q^2} + \frac{H}{2} = 0
\end{align}

Solving for \(Q\):
\begin{align}
\frac{D \cdot S}{Q^2} &= \frac{H}{2} \\
Q^2 &= \frac{2 \cdot D \cdot S}{H} \\
Q^* &= \sqrt{\frac{2 \cdot D \cdot S}{H}}
\end{align}

The second derivative confirms this is a minimum:
\begin{align}
\frac{d^2TC'(Q)}{dQ^2} &= \frac{2 \cdot D \cdot S}{Q^3} > 0 \quad \forall Q > 0
\end{align}

\begin{figure}[htbp]
\centering
\includegraphics[width=\textwidth]{../figures-png/line51.fig1.png}
\caption{Enhanced Economic Order Quantity Problem with Professional CMYK Colors and Node-Based Structure}
\end{figure}

\textbf{Concrete Example - MegaMart Coffee Beans:}
Using MegaMart's data:
- \(D = 12,000\) bags per year
- \(S = \$150\) per order
- \(H = \$8\) per bag per year

Calculating the optimal order quantity:
\begin{align}
Q^* &= \sqrt{\frac{2 \times 12,000 \times 150}{8}} \\
&= \sqrt{\frac{3,600,000}{8}} \\
&= \sqrt{450,000} \\
&= 671 \text{ bags}
\end{align}

Associated metrics:
- Number of orders per year: \(n^* = \frac{12,000}{671} = 17.9 \approx 18\) orders
- Order cycle length: \(\frac{365}{18} = 20.3\) days between orders
- Total annual cost: \(TC^* = \frac{12,000 \times 150}{671} + \frac{671 \times 8}{2} = 2,683 + 2,684 = \$5,367\)

\subsection{Tier 2: The Classic Heuristic (Python Implementation)}

The classical EOQ heuristic provides a straightforward computational approach to solving inventory optimization problems. This implementation focuses on the core EOQ calculation while incorporating practical considerations such as sensitivity analysis and constraint handling.

The algorithm implements the standard EOQ formula with additional features for practical application, including cost breakdown analysis, order scheduling, and parameter sensitivity evaluation. The approach uses a direct mathematical solution enhanced with practical business logic for real-world implementation.

\textbf{Algorithm: EOQ Classic Heuristic}

\textbf{Time Complexity:} \(O(1)\) for basic EOQ calculation, \(O(n)\) for sensitivity analysis with \(n\) parameter variations.

\textbf{Space Complexity:} \(O(1)\) for core calculation, \(O(n)\) for storing sensitivity results.

\begin{figure}[htbp]
\centering
\includegraphics[width=\textwidth]{../figures-png/line51.fig2.png}
\caption{EOQ Heuristic Algorithm Flowchart}
\end{figure}

\lstinputlisting[language=Python, caption=EOQ Classic Heuristic Implementation]{.././codes-python/line51.py1.py}

\textbf{Concrete Example Output:}

When executed with MegaMart's parameters, the EOQ heuristic produces the following results:

\begin{verbatim}
=== MegaMart Coffee Bean EOQ Analysis ===
Optimal Order Quantity: 670.82 bags
Orders per Year: 17.89
Days Between Orders: 20.4
Average Inventory Level: 335.41 bags
Annual Ordering Cost: $2,683.20
Annual Holding Cost: $2,683.28
Annual Inventory Cost: $5,366.48
Annual Purchasing Cost: $300,000.00
Total Annual Cost: $305,366.48

=== Sensitivity Analysis ===
Parameter Variation | Demand EOQ | Order Cost EOQ | Hold Cost EOQ
-----------------------------------------------------------------
       0.8          |   600.0    |     616.44     |     750.0
       0.9          |   636.4    |     653.2      |     707.1
       1.0          |   670.8    |     670.8      |     670.8
       1.1          |   703.3    |     687.13     |     639.12
       1.2          |   734.5    |     702.49     |     612.37

=== First Quarter Order Schedule ===
Day 1: Order 670.82 bags
Day 21: Order 670.82 bags
Day 42: Order 670.82 bags
Day 62: Order 670.82 bags
Day 83: Order 670.82 bags
\end{verbatim}

\subsection{Tier 3: The Advanced Algorithm (Dynamic EOQ with Demand Variability)}

The advanced EOQ algorithm extends the classic model to handle real-world complexities such as demand variability, lead times, and dynamic parameter updates. This implementation uses a Particle Swarm Optimization (PSO) approach to optimize EOQ parameters under uncertainty and multiple objectives.

The algorithm addresses limitations of the classic EOQ model by incorporating stochastic demand patterns, service level constraints, and multi-objective optimization. The PSO metaheuristic explores the solution space efficiently while balancing multiple competing objectives: cost minimization, service level maximization, and inventory turnover optimization.

\textbf{Algorithm: PSO-Enhanced Dynamic EOQ}

The Particle Swarm Optimization approach models each potential EOQ solution as a particle with position representing order quantity and reorder point decisions. Particles move through the solution space guided by personal best experiences and global swarm intelligence.

\textbf{Time Complexity:} \(O(g \times p \times f)\) where \(g\) is generations, \(p\) is population size, \(f\) is fitness evaluation complexity.

\textbf{Space Complexity:} \(O(p \times d)\) where \(p\) is population size, \(d\) is problem dimension.

\begin{figure}[htbp]
\centering
\includegraphics[width=\textwidth]{../figures-png/line51.fig3.png}
\caption{PSO-Enhanced Dynamic EOQ Algorithm Structure}
\end{figure}

\lstinputlisting[language=Python, caption=PSO-Enhanced Dynamic EOQ Implementation]{.././codes-python/line51.py2.py}

\textbf{Concrete Example Output:}

The PSO-enhanced dynamic EOQ optimization for MegaMart produces more robust results considering demand uncertainty:

\begin{verbatim}
=== MegaMart Dynamic EOQ with PSO Optimization ===
Optimizing dynamic EOQ with PSO...
Converged at iteration 67
Optimal Order Quantity: 748.32 bags
Optimal Reorder Point: 287.45 bags
Optimal Review Period: 18.7 days
Optimization Fitness: -2847.23
Convergence Iterations: 67
Recommended Safety Stock: 56.78 bags
\end{verbatim}

The advanced algorithm identifies that with demand variability, MegaMart should:
- Increase order quantity to 748 bags (vs. 671 in deterministic case)
- Maintain a reorder point of 287 bags to ensure 95\% service level
- Review inventory every 18.7 days for optimal responsiveness
- Hold approximately 57 bags of safety stock to buffer against demand uncertainty

The PSO approach successfully balances multiple objectives: minimizing total costs while maintaining high service levels and achieving efficient inventory turnover. The algorithm demonstrates superior performance in handling the complexity of stochastic demand patterns compared to deterministic EOQ models.

\section{Practice Questions}

\textbf{Tier 1 Questions (Mathematical Formulation):}

\textbf{Question 1:} A pharmaceutical distributor stocks a critical medication with the following parameters: annual demand of 8,400 units, ordering cost of \$200 per order, and holding cost of \$12 per unit per year. Formulate the complete mathematical model and calculate the optimal order quantity, number of orders per year, and total annual inventory cost. What is the order cycle length in days?

\textbf{Question 2:} A manufacturing company uses steel rods with annual demand of 15,000 units, ordering cost of \$300, and holding cost varying seasonally between \$5-\$15 per unit per year. Develop a mathematical sensitivity analysis showing how EOQ changes with holding cost variations. At what holding cost does the EOQ equal 1,000 units?

\textbf{Tier 2 Questions (Classic Heuristic):}

\textbf{Question 3:} Implement the EOQ heuristic for a bookstore chain with the following data: 25 locations, each selling an average of 480 copies per year of a bestselling novel, combined ordering cost of \$75 per order, holding cost of \$3.50 per book per year, and book cost of \$18. Calculate the optimal centralized ordering strategy and compare total costs with decentralized ordering (each store orders independently). Include a sensitivity analysis for ±20\% demand variations.

\textbf{Question 4:} A grocery chain manages frozen food inventory with daily demand following a normal distribution (mean = 45 units, standard deviation = 12 units), ordering cost of \$125, holding cost of \$2.80 per unit per year, and a 3-day lead time. Extend the basic EOQ heuristic to incorporate safety stock calculations for 90\%, 95\%, and 99\% service levels. What is the impact on total costs for each service level?

\textbf{Tier 3 Questions (Advanced Metaheuristic):}

\textbf{Question 5:} Design a multi-objective EOQ optimization for an electronics retailer managing 50 different smartphone models. Each model has different demand patterns (provided dataset includes monthly sales for 2 years), ordering costs ranging from \$80-\$500, holding costs varying by model value (\$15-\$200 per unit per year), and obsolescence risks (3-18 month product lifecycles). Use a metaheuristic approach to optimize across multiple objectives: cost minimization, service level maximization (target: 98\%), and inventory turnover maximization (target: 8 turns per year). Include constraint handling for warehouse capacity (10,000 units total) and cash flow limits (\$2M maximum inventory investment).

\textbf{Question 6:} Develop an adaptive EOQ system using evolutionary algorithms for a seasonal fashion retailer. The system must handle: (1) Demand patterns with 300\% seasonal variation between peak and off-peak periods, (2) Dynamic pricing strategies affecting holding costs, (3) Supplier capacity constraints limiting maximum order quantities, (4) Multi-echelon inventory considerations with 3 distribution centers and 150 retail locations, and (5) Markdown optimization for end-of-season inventory. Implement and test the algorithm with provided historical data showing demand seasonality, price elasticity parameters, and supplier constraints. Compare performance against traditional EOQ approaches across multiple performance metrics including total cost, service level, inventory turnover, and markdown losses.

\end{document}
